\FloatBarrier
\section{本章小结}
\label{sec:matrix-equations-chapter-summary}

本章研究以矩阵为未知量、同时受左乘和右乘作用的方程。左右两种作用相加以后不能分别“除以”系数，但左边仍是矩阵空间上的线性算子；向量化把它写成Kronecker系数矩阵，说明一般线性方程的判据都适用，却不宜作为大规模计算的默认方法。

Sylvester方程$\bm A\bm X+\bm X\bm B=\bm C$对任意右端唯一可解，当且仅当$\bm A$与$-\bm B$没有公共特征值，这与两侧系数各自是否可逆无关。谱和为零时，对角情形清楚地显示出无解与多解两种退化。谱分离只回答唯一性，不回答敏感度：最小谱和接近零，或者特征基远离正交，都会放大数据误差。

计算上，Bartels--Stewart算法用两个正交Schur变换把方程化为准三角形，按列从左到右、每列从下到上回代，二阶对角块只形成至多四个未知数的局部系统；正交变换保持Frobenius残差，所以原方程的残差可以直接检验整个流程。可对角化时换到特征基可以逐元素相除，但只有特征基接近正交时才可靠，一般情形仍用Schur路线。

Lyapunov方程来自二次量的累计与变化。离散型$\bm X-\bm A^\top\bm X\bm A=\bm Q$唯一可解的条件是特征值两两乘积不为一，连续型$\bm A^\top\bm X+\bm X\bm A=-\bm Q$的条件是两两之和不为零；右端对称时解对称。在本章的写法下，两类方程都只需对$\bm A$作一次Schur分解，再按列从左到右、每列从上到对角元逐元素递推，只求上三角部分。浮点结果要在原方程中检查残差，对称化和Cholesky分解分别回答对称性与正定性，不能互相代替。

正定性方面，$\rho(\bm A)<1$时离散型的解就是绝对收敛的级数$\sum_k(\bm A^\top)^k\bm Q\bm A^k$：右端正定时解正定；右端只半正定时，解正定当且仅当不存在始终不被累计量察觉的非零初始方向，这一条件只需检查$\bm Q\bm A^k\bm x=\bm0$（$k<d$）。反过来，对称正定解本身就是谱的证书：离散型给出$\rho(\bm A)<1$，连续型给出特征值实部全负，证明只用代数运算。

连续型方程在特征值实部全负时的正定解，本章只陈述了条件并用例子与谱证书作代数检验。它的一般证明要把离散级数换成矩阵指数的无穷积分，需要线性微分方程的解理论；这一积分表示连同它与轨迹长期稳定性的关系，在后续关于演化的理论中建立（见第~\ref{sec:ls-lyapunov-positivity}节所指的位置）。
